% GENERATED FILE. Edit canonical lessons, then run npm run generate:books.
% canonical-source-hash: fnv1a64:b0725b9b45bb3bcd
% canonical-lessons: UR-C340-faisla, UR-C340-zamanat, UR-C340-mulzim, UR-C340-munsif, UR-C340-gavahi-dena

\chapter{In Court}
\label{ch:ur-in-court}
\hlchaptermodality{340}

\begin{chapteropening}
\textbf{By the end of this chapter:} \emph{I can follow what happens in a court case.}

Complete the last lesson of chapter 340: I can follow what happens in a court case.
\end{chapteropening}

\section[\texorpdfstring{\ur{فیصلہ} (faislā)}{فیصلہ (faislā)}]{\ur{فیصلہ} (faislā) — a verdict; a decision}

\label{lesson:UR-C340-faisla}



\begin{quote}
Before the new one: say the Urdu for to act on, then the Urdu for to fulfil.
\end{quote}

\subsection*{You'll want to know: \ur{فیصلہ}}
\textbf{\ur{فیصلہ}} — \emph{faislā} — \textquotedblleft{}a verdict; a decision\textquotedblright{}.

\begin{quote}
\ur{جج} \ur{نے} \ur{آج} \ur{فیصلہ} \ur{دیا۔}
\end{quote}

\begin{quote}
\emph{jaj ne āj faislā diyā.}
\end{quote}

\begin{quote}
The judge gave the verdict today.
\end{quote}

Masculine.

\begin{cousinweb}
\textbf{\ur{فیصلہ}} (\emph{faislā}) is usually traced, through Persian, to Arabic \emph{faiṣala}, \textquotedblleft{}a decision\textquotedblright{}.
\end{cousinweb}

\begin{grammarlens}[title={What you've built}]
One more word for giving reasons and opinions.
\end{grammarlens}

\subsection*{Your turn}
\begin{itemize}
\raggedright
  \item \emph{Say it:} \emph{faislā}
  \item \emph{Say it:} \emph{faislā}, once more
  \item \emph{Say it:} \emph{jaj ne āj faislā diyā.} — The judge gave the verdict today.
  \item \emph{Recall:} say the Urdu for to act on, then the Urdu for to fulfil, then say \emph{faislā} again
\end{itemize}

\begin{tcolorbox}[breakable,colback=teal!4,colframe=teal!35!black,title={Before you move on}]
What does \ur{فیصلہ} mean? (\textquotedblleft{}A verdict; a decision\textquotedblright{}.) Say it once more.
\end{tcolorbox}
\section[\texorpdfstring{\ur{ضمانت} (zamānat)}{ضمانت (zamānat)}]{\ur{ضمانت} (zamānat) — bail; a guarantee}

\label{lesson:UR-C340-zamanat}



\begin{quote}
Before the new one: say the Urdu for to fulfil, then the Urdu for a verdict.
\end{quote}

\subsection*{You'll want to know: \ur{ضمانت}}
\textbf{\ur{ضمانت}} — \emph{zamānat} — \textquotedblleft{}bail; a guarantee\textquotedblright{}.

\begin{quote}
\ur{اسے} \ur{ضمانت} \ur{مل} \ur{گئی۔}
\end{quote}

\begin{quote}
\emph{use zamānat mil gaī.}
\end{quote}

\begin{quote}
He got bail.
\end{quote}

Feminine.

\begin{cousinweb}
\textbf{\ur{ضمانت}} (\emph{zamānat}) is usually traced, through Persian, to Arabic \emph{ḍamāna}, \textquotedblleft{}surety\textquotedblright{}: a pledge that someone will come back.
\end{cousinweb}

\begin{grammarlens}[title={What you've built}]
One more word for giving reasons and opinions.
\end{grammarlens}

\subsection*{Your turn}
\begin{itemize}
\raggedright
  \item \emph{Say it:} \emph{zamānat}
  \item \emph{Say it:} \emph{zamānat}, once more
  \item \emph{Say it:} \emph{use zamānat mil gaī.} — He got bail.
  \item \emph{Recall:} say the Urdu for to fulfil, then the Urdu for a verdict, then say \emph{zamānat} again
\end{itemize}

\begin{tcolorbox}[breakable,colback=teal!4,colframe=teal!35!black,title={Before you move on}]
What does \ur{ضمانت} mean? (\textquotedblleft{}Bail; a guarantee\textquotedblright{}.) Say it once more.
\end{tcolorbox}
\section[\texorpdfstring{\ur{ملزم} (mulzim)}{ملزم (mulzim)}]{\ur{ملزم} (mulzim) — an accused person}

\label{lesson:UR-C340-mulzim}



\begin{quote}
Before the new one: say the Urdu for a verdict, then the Urdu for bail.
\end{quote}

\subsection*{You'll want to know: \ur{ملزم}}
\textbf{\ur{ملزم}} — \emph{mulzim} — \textquotedblleft{}an accused person\textquotedblright{}.

\begin{quote}
\ur{ملزم} \ur{نے} \ur{جج} \ur{سے} \ur{کچھ} \ur{نہیں} \ur{کہا۔}
\end{quote}

\begin{quote}
\emph{mulzim ne jaj se kuchh nahīṅ kahā.}
\end{quote}

\begin{quote}
The accused said nothing to the judge.
\end{quote}

Masculine or feminine, by the person meant.

\begin{cousinweb}
\textbf{\ur{ملزم}} (\emph{mulzim}) is usually traced, through Persian, to Arabic \emph{mulzam}, \textquotedblleft{}one held liable\textquotedblright{}; Urdu says it \emph{mulzim} and uses it for a person charged with a crime. It is of the same family of words as \textbf{\ur{الزام}}, \textquotedblleft{}an accusation\textquotedblright{}.
\end{cousinweb}

\begin{grammarlens}[title={What you've built}]
One more word for giving reasons and opinions.
\end{grammarlens}

\subsection*{Your turn}
\begin{itemize}
\raggedright
  \item \emph{Say it:} \emph{mulzim}
  \item \emph{Say it:} \emph{mulzim}, once more
  \item \emph{Say it:} \emph{mulzim ne jaj se kuchh nahīṅ kahā.} — The accused said nothing to the judge.
  \item \emph{Recall:} say the Urdu for a verdict, then the Urdu for bail, then say \emph{mulzim} again
\end{itemize}

\begin{tcolorbox}[breakable,colback=teal!4,colframe=teal!35!black,title={Before you move on}]
What does \ur{ملزم} mean? (\textquotedblleft{}An accused person\textquotedblright{}.) Say it once more.
\end{tcolorbox}
\section[\texorpdfstring{\ur{منصف} (munsif)}{منصف (munsif)}]{\ur{منصف} (munsif) — an arbiter; a fair judge}

\label{lesson:UR-C340-munsif}



\begin{quote}
Before the new one: say the Urdu for bail, then the Urdu for an accused person.
\end{quote}

\subsection*{You'll want to know: \ur{منصف}}
\textbf{\ur{منصف}} — \emph{munsif} — \textquotedblleft{}an arbiter; a fair judge\textquotedblright{}.

\begin{quote}
\ur{اچھا} \ur{منصف} \ur{دونوں} \ur{فریقوں} \ur{کی} \ur{بات} \ur{سنتا} \ur{ہے۔}
\end{quote}

\begin{quote}
\emph{achchhā munsif donoṅ farīqoṅ kī bāt suntā hai.}
\end{quote}

\begin{quote}
A good arbiter listens to both sides.
\end{quote}

Masculine or feminine, by the person meant.

\begin{cousinweb}
\textbf{\ur{منصف}} (\emph{munsif}) is usually traced, through Persian, to Arabic \emph{munṣif}, \textquotedblleft{}one who is fair\textquotedblright{}, of the same family as \textbf{\ur{انصاف}}, \textquotedblleft{}justice\textquotedblright{}.
\end{cousinweb}

\begin{grammarlens}[title={What you've built}]
One more word for giving reasons and opinions.
\end{grammarlens}

\subsection*{Your turn}
\begin{itemize}
\raggedright
  \item \emph{Say it:} \emph{munsif}
  \item \emph{Say it:} \emph{munsif}, once more
  \item \emph{Say it:} \emph{achchhā munsif donoṅ farīqoṅ kī bāt suntā hai.} — A good arbiter listens to both sides.
  \item \emph{Recall:} say the Urdu for bail, then the Urdu for an accused person, then say \emph{munsif} again
\end{itemize}

\begin{tcolorbox}[breakable,colback=teal!4,colframe=teal!35!black,title={Before you move on}]
What does \ur{منصف} mean? (\textquotedblleft{}An arbiter; a fair judge\textquotedblright{}.) Say it once more.
\end{tcolorbox}
\section[\texorpdfstring{\ur{گواہی} \ur{دینا} (gavāhī denā)}{گواہی دینا (gavāhī denā)}]{\ur{گواہی} \ur{دینا} (gavāhī denā) — to testify; to bear witness}

\label{lesson:UR-C340-gavahi-dena}



\begin{quote}
Before the new one: say the Urdu for an accused person, then the Urdu for an arbiter.
\end{quote}

\subsection*{You'll want to know: \ur{گواہی} \ur{دینا}}
\textbf{\ur{گواہی} \ur{دینا}} — \emph{gavāhī denā} — \textquotedblleft{}to testify; to bear witness\textquotedblright{}.

\begin{quote}
\ur{سارا} \ur{نے} \ur{اس} \ur{کے} \ur{حق} \ur{میں} \ur{گواہی} \ur{دی۔}
\end{quote}

\begin{quote}
\emph{sārā ne us ke haq meṅ gavāhī dī.}
\end{quote}

\begin{quote}
Sara testified in his favour.
\end{quote}

Past, with \ur{نے} (he or she …): \ur{اس} \ur{نے} \ur{گواہی} \ur{دی}, \emph{us ne gavāhī dī}.

\begin{cousinweb}
\textbf{\ur{گواہی} \ur{دینا}} (\emph{gavāhī denā}) is built from \emph{gavāhī}, \textquotedblleft{}testimony\textquotedblright{}, usually traced to Persian \emph{gavāhī}, from \emph{gavāh}, \textquotedblleft{}a witness\textquotedblright{}, and \textbf{\ur{دینا}}, \textquotedblleft{}to give\textquotedblright{}.
\end{cousinweb}

\begin{grammarlens}[title={What you've built}]
One more word for giving reasons and opinions.
\end{grammarlens}

\subsection*{Your turn}
\begin{itemize}
\raggedright
  \item \emph{Say it:} \emph{gavāhī denā}
  \item \emph{Say it:} \emph{gavāhī denā}, once more
  \item \emph{Say it:} \emph{sārā ne us ke haq meṅ gavāhī dī.} — Sara testified in his favour.
  \item \emph{Recall:} say the Urdu for an accused person, then the Urdu for an arbiter, then say \emph{gavāhī denā} again
\end{itemize}

\begin{tcolorbox}[breakable,colback=teal!4,colframe=teal!35!black,title={Before you move on}]
What does \ur{گواہی} \ur{دینا} mean? (\textquotedblleft{}To testify; to bear witness\textquotedblright{}.) Say it once more.
\end{tcolorbox}
